% THIS IS A LATEX TEMPLATE FILE FOR PAPERS INCLUDED IN THE
% *Anthology of Computers and the Humanities*.
\documentclass[final]{anthology-ch}

% LOAD LaTeX PACKAGES
\usepackage{booktabs}
\usepackage{graphicx}

% TITLE OF THE SUBMISSION
\title{More Than Just Planting Pansies: Black Women's Garden Clubs and the Neo-Archive in North Carolina's Digital Landscape}

% AUTHOR AND AFFILIATION INFORMATION
\author[1]{Rhonda D. Jones}[
  orcid=
]

\author[2]{Maurika Smutherman}[
  orcid=
]

\affiliation{1}{University of North Carolina at Greensboro, Greensboro, NC, USA}
\affiliation{2}{North Carolina Central University, Durham, NC, USA}

% KEYWORDS
\keywords[eng]{African American women, digital public history, Black environmental stewardship, community archives, oral history}

% METADATA FOR THE PUBLICATION
\pubyear{2025}
\pubvolume{5}
\pagestart{77}
\pageend{83}
\paperorder{10}
\conferencename{From Manuscript to Megabyte: Building Sustainable and Open Futures in the Digital Humanities}
\conferenceeditors{James B. Harr III, Emma Stanley, Jenna Rinalducci, and Donna Kain}
\doi{10.63744/dVeLHNjFxvtP}
\customhead{}

\addbibresource{bibliography.bib}

%%%%%%%%%%%%%%%%%%%%%%%%%%%%%%%%%%%%%%%%%%%%%%%%%%%%%%%%%%%%%%%%%%%%%%%%%%%
% HERE IS THE START OF THE TEXT
\begin{document}

\maketitle

\begin{abstract}
African American women's garden clubs played a significant yet understudied role in the civic, environmental, and cultural life of Black communities throughout the Jim Crow South. While historians have documented the importance of Black women's club movements in education, mutual aid, and political activism, garden clubs have often been reduced to organizations concerned primarily with beautification and leisure. This article argues that Black women's garden clubs functioned as institutions of environmental stewardship, civic leadership, historical preservation, and community care. Recovering these histories presents a methodological challenge because documentation survives not within centralized repositories but across newspapers, scrapbooks, oral histories, family collections, photographs, and community memory.

Drawing on digital humanities research into the North Carolina Federation of Negro Garden Clubs, the leadership of Mrs. Madie Hall Xuma, local garden club records, oral histories, and community-generated collections, this article examines how digital methods can reconstruct fragmented histories while revealing the limitations of traditional archival systems. Using newspapers, oral history, relational metadata, spatial mapping, and Omeka S, the project proposes the concept of the neo-archive, a framework that expands understandings of historical evidence by recognizing gardens, landscapes, beauty practices, oral knowledge, and memory as legitimate archival forms. In doing so, it contributes to scholarship on Black women's civic leadership, Black environmental history, community archives, and how digital humanities can recover overlooked histories while challenging assumptions about where historical knowledge resides and how it is preserved.
\end{abstract}

\section{Introduction} \label{sec:intro}

In 1971, the \textit{Durham Morning Herald} profiled Audrey Meadows and the garden she inherited from her mother, Mrs.\ Haddie Meadows. The article emphasized that successful gardening required \enquote{work and planning,} describing a carefully maintained landscape shaped through years of labor, environmental knowledge, and intergenerational care~\cite{hodges1971green}. At first glance, the story appears to be a routine feature about gardening. Yet for historians seeking to reconstruct the history of African American women's garden clubs, the article reveals something more significant: the transmission of ecological knowledge across generations, the cultivation of family memory through landscape, and the enduring legacy of Black women's community leadership.

The article also exposes a larger archival problem. How do historians recover the histories of Black women's civic institutions when their records remain fragmented, dispersed, and only partially preserved? Although African American women established extensive civic organizations throughout the twentieth century, the documentary record of Black women's garden clubs remains uneven and scattered. Unlike many political organizations and professional associations, garden clubs rarely deposited their records within a single institutional repository. Instead, their histories survive across local newspapers, family scrapbooks, photographs, oral histories, newsletters, websites, blogs, and social media accounts. The absence of a centralized archive has contributed to the misconception that these organizations occupied only a marginal place in Black civic life.

This article argues the opposite. Black women's garden clubs were obscured not by historical insignificance but by archival fragmentation. Recovering their histories requires historians to reconsider what constitutes historical evidence and to develop methodologies capable of reconstructing organizations whose records remain dispersed across multiple forms of community documentation. Rather than treating newspapers, scrapbooks, landscapes, and oral histories as separate categories of evidence, this article reads them relationally, as parts of a distributed archive.

This argument builds upon several strands of scholarship. The rich literary scholarship of Alice Walker, Maya Angelou and historiography of Black women's club movements, Elsa Barkley Brown, Angela Davis, Jacqueline Jones, Cynthia Neverdone-Morton, Stephanie Shaw, Deborah Gray White, and Glenda Gilmore, have demonstrated that African American women created extensive networks of civic organizations dedicated to Progressive Era reform, education, mutual aid, racial uplift, and community improvement. Their work established Black women's clubs as important institutions of political and social leadership during segregation. Yet within this rich historiography, garden clubs have received comparatively little sustained scholarly attention, often appearing only briefly within broader studies of clubwomen's activism.

Black environmental historians have begun to fill this gap by reframing gardening as a form of environmental knowledge and civic practice rather than merely domestic labor or aesthetic expression. Dianne Glave argues that African American women gardeners cultivated landscapes that reflected ecological knowledge and community improvement, expanding traditional understandings of environmental history beyond conservation organizations and wilderness preservation~\cite{glave2003garden}. Glave and Mark Stoll further situate African American relationships with land within histories of stewardship, community, and cultural survival, while Abra Lee documents the overlooked contributions of Black gardeners, horticulturalists, and landscape professionals to American horticultural history~\cite{glave2006love, lee2021conquer}. Together, this scholarship demonstrates that gardening constituted an important dimension of African American environmental thought, but it has not fully examined how garden clubs histories might be reconstructed through digital methods.

The life of Mrs.\ Madie Hall Xuma illustrates these intersections particularly well. Historian Wanda Hendricks demonstrates that Xuma's leadership in founding Winston-Salem's \enquote{Garden Path Flower Club} and serving as the first president of the North Carolina Federation of Negro Garden Clubs preceded her internationally recognized activism in South Africa as the first president of the African National Congress Women's League~\cite{hendricks2022life}. Xuma's career reveals that local garden clubs were not peripheral social organizations but important sites of leadership development whose influence extended far beyond horticulture.

Building upon these conversations, I propose the concept of the \textit{neo-archive}. The neo-archive is not an alternative repository but a methodological framework for recognizing and connecting distributed forms of historical evidence. It acknowledges that African American communities have long preserved historical knowledge through garden landscapes, seed cultivation, beautification practices, photographs, and scrapbooks. Rather than treating these materials as supplementary to institutional archives, the neo-archive recognizes them as interconnected archival systems capable of documenting histories that traditional repositories only partially capture.

Combining newspaper research, community collaboration, oral histories, relational metadata, and the development of a digital archive in Omeka S, this article contributes to three intersecting fields of scholarship. First, it expands scholarship on Black women's club movements by examining garden clubs as important sites of leadership and civic engagement. Second, it contributes to Black environmental history by exploring how African American women cultivated landscapes that reflected environmental knowledge, cultural values, and community aspirations. Third, it intervenes in digital humanities and archival studies by demonstrating how digital methods can reconstruct histories that remain fragmented across distributed sources.

The study reconstructs the history of Black women's garden clubs by linking dispersed forms of evidence into a coherent historical narrative. In doing so, it demonstrates that digital humanities can serve not only as a means of expanding access to archival materials but also as a methodology for historical recovery, one capable of revealing relationships among people, places, and practices that have long remained obscured.

\section{Black Women's Garden Clubs in North Carolina: Civic Leadership, Environmental Stewardship, and Community Memory} \label{sec:clubs}

Black women's garden clubs developed during a period when African Americans were systematically excluded from many of the public spaces, civic institutions, and environmental amenities that shaped everyday life in the Jim Crow South. Parks, recreational facilities, municipal beautification projects, and neighborhood improvements frequently reflected racial inequities in public investment. In response, African American communities created their own institutions through churches, schools, mutual aid societies, women's clubs, and neighborhood organizations. Garden clubs emerged within this broader tradition of institution-building, transforming horticulture into a form of civic engagement through which Black women cultivated not only flowers and landscapes but also leadership, education, and community identity.

Unlike many white garden clubs that concentrated primarily on ornamental horticulture, African American garden clubs often linked beautification with racial uplift, public health, education, and neighborhood improvement. Meetings combined horticultural instruction with discussions of civic responsibility, while flower shows, home-garden competitions, and beautification campaigns encouraged residents to improve both private and public spaces. Gardening became a visible expression of dignity and self-determination within communities that had long been denied equal access to public resources. As Glave argues, African American women's gardening practices reflected environmental knowledge rooted in community improvement and social reform rather than simply aesthetic preference~\cite[pp.~399--404]{glave2003garden}.

North Carolina became an important center for this movement. Throughout the mid-twentieth century, local clubs organized educational programs, exchanged horticultural knowledge, sponsored community improvement projects, and participated in regional and statewide networks. Although these organizations rarely left extensive institutional archives, newspaper coverage demonstrates that they were remarkably active. Reports in \textit{The Carolina Times}, \textit{The Durham Morning Herald}, and \textit{Winston Salem Journal} documented regular meetings, officer elections, flower shows, youth gardening initiatives, and federation conventions across the state. When read collectively rather than as isolated announcements, these articles reveal a sustained infrastructure of Black women's environmental leadership extending over several decades.

Among the most influential figures in this movement was Mrs.\ Madie Hall Xuma, whose career illustrates the relationship between local civic engagement and global Black women's leadership. In 1931, Xuma an avid gardener, organized the \enquote{Garden Path Flower Club} in Winston-Salem, recognized as the city's first African American garden club. Four years later, she helped establish the North Carolina Federation of Negro Garden Clubs at North Carolina Agricultural and Technical College, serving as its first president. The use of vegetable-plant gardens culminated to combating hunger and poverty, also transformed littered strewn yards into beautified spaces. Under her leadership, the federation united local clubs across the state through annual conventions, educational programming, horticultural competitions, and leadership development~\cite[p.~72]{hendricks2022life}.

Hendricks situates Xuma's garden club leadership within a much broader history of Black women's activism. Rather than viewing her horticultural work as separate from her later political career, Hendricks demonstrates that the organizational skills Xuma developed in North Carolina became foundational to her subsequent leadership in South Africa, where she served as the first president of the African National Congress Women's League~\cite[p.~136]{hendricks2022life}. Seen from this perspective, the North Carolina Federation of Negro Garden Clubs was more than a horticultural organization; it was an institution through which Black women developed administrative expertise, public speaking skills, regional networks, and models of collaborative leadership.

The federation also challenged conventional distinctions between environmental work and political engagement. Beautification projects were not merely decorative exercises but investments in community well-being. Club members planted flowers around churches, schools, cemeteries, and neighborhood gathering spaces, organized educational lectures on horticulture and conservation, and encouraged home gardening as both an aesthetic and practical endeavor. These activities reflected what Glave and Stoll describe as an African American environmental tradition grounded in stewardship, community, and cultural survival~\cite{glave2006love}. Garden clubs cultivated landscapes that were simultaneously ecological, educational, and social.

The significance of this work becomes particularly evident when examined through local examples. The \enquote{Lyon Park Rose Garden Club} in Durham illustrates how neighborhood organizations connected environmental stewardship with civic life. Newspaper accounts describe regular meetings, demonstrations, juried competitions, seasonal flower displays, and participation in wider federation activities. Yet these published accounts represent only one layer of the historical record. Photographs preserved by club members and neighborhood landscapes reveal additional dimensions of club life that rarely appeared in formal reports.

The \textit{Les Fleurs Garden Club Scrapbook}, preserved by the Durham County Library Digital Archives, offers an especially important example of this distributed historical record. The scrapbook contains photographs, newspaper clippings, invitations, programs, handwritten annotations, and other ephemera documenting decades of organizational activity~\cite{lesfleurs}. Rather than functioning simply as a commemorative album, it demonstrates how club members consciously documented their own history. The scrapbook records club members, public events, floral exhibitions, and civic accomplishments that remain largely absent from institutional archives. Read alongside newspapers and oral histories, it reveals how Black women created their own mechanisms of historical preservation even when their organizations received limited attention from collecting institutions.

Taken together, these sources complicate traditional understandings of historical evidence. Newspapers document organizational activity but often omit personal relationships and community meaning. Scrapbooks preserve visual and social memory while revealing how members wished to represent themselves. Oral histories provide context unavailable within written records, explaining why particular gardens, neighborhoods, and civic projects mattered to participants. Landscapes themselves preserve traces of environmental knowledge through plantings, spatial organization, and long-standing gardening traditions. No single source is sufficient. Instead, historical meaning emerges through the relationships among them.

\section{The Archival Problem, the Neo-Archive, and Digital Humanities Methodology} \label{sec:method}

The history of Black women's garden clubs presents historians with a paradox. Newspaper databases contain scattered references to meetings, flower shows, conventions, officer elections, and community beautification projects, yet no single archival collection documents the movement in its entirety. Scrapbooks preserve photographs and programs but rarely explain the organizational networks to which they belonged. Oral histories illuminate community relationships but often lack precise dates or documentary corroboration. Family collections preserve intergeneration relationships while remaining largely inaccessible to researchers. Individually, these sources appear incomplete. Collectively, however, they document a sophisticated network of Black women's environmental leadership.

This fragmentation is not unique to garden clubs. Traditional archives have long privileged institutional records, governmental documents, and organizational papers created by those with the resources to preserve them. Community organizations, particularly those led by African American women, often produced records that remained within homes, churches, neighborhood organizations, and personal collections rather than entering university or state archives. As a result, historians encounter not the absence of history but the absence of centralized documentation.

The research process for the Black Women's Garden Clubs of North Carolina project made this reality immediately apparent. Initial searches sought organizational records comparable to those available for many civic associations. Instead, evidence surfaced incrementally through digitized Black newspapers, county historical collections, family scrapbooks, oral histories, cemetery records, community websites, and local archival repositories. Each source documented only a portion of the historical record. The challenge therefore became not locating a single archive but identifying relationships among dispersed sources.

Digital humanities provide an important methodological response to this problem. Rather than viewing archival fragments as isolated documents, digital methods make it possible to model relationships among people, organizations, places, events, and material culture. Within the \textit{More Than Just Planting Pansies Neo-Archives} project, Omeka S serves not simply as a repository for digitized materials but as a relational knowledge environment. Individual items, including newspaper articles, photographs, convention programs, club histories, oral history excerpts, and scrapbook pages, are connected through shared metadata describing people, organizations, geographic locations, events, and themes. These relationships allow researchers to reconstruct organizational networks that are not visible when sources are examined independently.

This relational approach differs from traditional digitization projects, which frequently reproduce the structure of existing archival collections. If archival silences remain unexamined, digitization can simply replicate historical inequities in digital form. The goal here is not merely to digitize surviving materials but to recover histories fragmented by unequal collecting practices and institutional neglect.

The concept of the neo-archive emerged directly from this research process. Rather than defining the archive solely as a repository of documents held within institutional custody, the neo-archive recognizes historical knowledge as distributed across community-generated forms of evidence. It acknowledges that Black communities have long preserved history through practices of stewardship, memory, and collective care that fall outside conventional archival categories. Newspapers, gardens, oral traditions, family photographs, scrapbooks, beautification practices, and neighborhood landscapes are therefore understood not as supplementary evidence but as interconnected archival systems.

This expanded understanding of evidence has important implications for the history of Black women's garden clubs. Seed-saving, for example, documents the transmission of environmental knowledge across generations. Landscape design reflects aesthetic values, ecological expertise, and community aspirations. Beautification projects record investments in public space, neighborhood identity, and civic responsibility. Oral histories preserve relationships, organizational memory, and local knowledge unavailable within written sources. Community scrapbooks curate historical narratives created by participants themselves rather than institutional collectors. Together, these materials reveal histories that remain largely invisible within conventional archival repositories.

As previously mentioned, the \textit{Les Fleurs Garden Club Scrapbook} illustrates this process particularly well~\cite{lesfleurs}. The commemorative album functions as an intentional act of community documentation. Whereas participants selected what should be remembered, preserved visual evidence of organizational achievements, and recorded relationships that would otherwise disappear from the historical record. When interpreted alongside newspaper coverage and oral testimony, the scrapbook becomes both a historical source and an archival intervention: evidence that Black women actively documented their own communities even when formal collecting institutions did not.

Community collaboration further expands this methodological framework. Through partnerships with descendants, local historians, and community organizations, the project recognizes that historical expertise does not reside exclusively within academic institutions. Community members frequently possess photographs, family histories, club memorabilia, and contextual knowledge unavailable in archival repositories. Their participation not only enriches the historical record but also redistributes authority within the research process. Digital humanities become a collaborative practice of historical reconstruction rather than simply a technological exercise.

This methodology is also informed by trauma-informed archival practice. The fragmented archival record of Black women's organizations reflects histories of racial segregation, unequal institutional investment, and selective preservation rather than historical insignificance. Approaching these absences through a trauma-informed framework shifts the research question from \textit{Why are the records missing?} to \textit{How did communities preserve knowledge despite archival exclusion?} This perspective recognizes resilience alongside loss and foregrounds the creative strategies through which Black communities maintained historical continuity.

The neo-archive therefore represents more than a new category of evidence; it offers a methodological intervention within digital humanities. It challenges scholars to move beyond document-centered models of archival practice by recognizing that historical knowledge is often embodied, relational, and community-generated. For Black women's garden clubs, history survives not within a single repository but across interconnected networks of newspapers, photographs, oral histories, landscapes, scrapbooks, and living memory. Digital humanities provide the tools to reconnect these fragments, while the neo-archive provides the theoretical framework for understanding why those connections matter.

\section{Conclusion} \label{sec:conclusion}

The history of Black women's garden clubs in North Carolina demonstrates that some of the most significant stories of African American civic life survive not within centralized archival repositories but across dispersed networks of community memory, newspapers, photographs, scrapbooks, oral histories, and living landscapes. These fragmented records have contributed to the mistaken perception that garden clubs occupied only a marginal place within African American history. Read collectively, however, they reveal a sophisticated network of organizations that cultivated environmental stewardship, neighborhood improvement, civic leadership, and intergenerational knowledge during the Jim Crow era.

This article has argued that Black women's garden clubs should be understood as institutions of environmental citizenship. Through organizations such as the North Carolina Federation of Negro Garden Clubs and local clubs, African American women transformed gardening into a form of public leadership that linked ecological knowledge with education, beautification, and community care. Their work challenges conventional distinctions between aesthetics and politics, demonstrating that the cultivation of gardens also constituted the cultivation of communities.

Recovering this history requires historians to reconsider the nature of archives themselves. Traditional repositories remain indispensable to historical scholarship, yet they cannot fully account for communities whose histories were preserved through distributed and often informal systems of documentation. Local newspaper coverage, family collections, oral histories, neighborhood landscapes, artifacts, and ephemera each preserve distinct dimensions of the historical record. Individually they appear incomplete; together they document organizational networks, environmental practices, and community relationships that would otherwise remain obscured.

The concept of the neo-archive offers a framework for understanding these relationships. Rather than replacing conventional archives, the neo-archive expands the evidentiary landscape by recognizing that historical knowledge is also preserved through landscapes, beauty practices, oral traditions, seed-sharing, family stewardship, and community memory. These forms of preservation have long sustained African American historical knowledge despite uneven collecting practices and the relative absence of institutional documentation. Recognizing them as archival evidence broadens not only the historical record but also our understanding of how communities intentionally preserve their own histories.

The project also demonstrates how digital humanities can move beyond digitization toward historical reconstruction. Omeka S, relational metadata, newspapers, and community collaboration do more than increase access to historical materials; they reconnect dispersed evidence into meaningful historical relationships. Digital humanities therefore become a methodology for recovering histories fragmented by archival silences rather than simply a platform for presenting existing collections. By linking newspapers to photographs, oral histories to landscapes, and scrapbook pages to organizational records, digital tools reveal networks of people, places, and practices that remain largely invisible when sources are examined in isolation.

Ultimately, Black women's garden clubs compel us to reconsider where history resides. Their legacy is preserved not only in institutional repositories but also in gardens that continue to flourish, in photographs carefully preserved by families, in scrapbooks assembled by club members, in newspaper columns chronicling neighborhood life, and in stories passed from one generation to the next. These are not peripheral sources to be consulted only when official records are absent. They are archives in their own right: repositories of environmental knowledge, civic leadership, and collective memory that expand both the practice of history and the possibilities of digital humanities.

\printbibliography

\end{document}
